\section{The socle on 120 Fano subplanes}
Use the fixed $S(4,7,23)$ already constructed and checked by
\texttt{steiner\_construction.py}. Delete two named points. The $120$ heptads
avoiding them form a set $\mathcal H$. The $21$ heptads containing both
points, with those two removed, form the lines of a projective plane on
$21$ points. Its lines have size five; every two remaining points lie on
one line by the Steiner property. Two such lines meet in one point, since
the original heptads meet in one or three points and already share the
deleted pair.

For $H\in\mathcal H$, every two points of $H$ lie on a plane line meeting
$H$ in three points. Indeed the heptad defining that line shares at least
two points with $H$ and hence exactly three. These triples form the Fano
subplane on the seven points of $H$. Let $\mathcal F$ be the $105$ incident
point--line flags of the plane, and define a $105\times120$ matrix $R$ by
\[
 R_{(p,\ell),H}=1\quad\Longleftrightarrow\quad
 p\in H\ \hbox{ and }\ |\ell\cap H|=3.
\]
Thus each column records the $21$ flags of the corresponding Fano subplane.
All following code spaces and matrices are over $\mathbb F_3$.

The common square-zero incidence lemma \cite[Lemma~6.1]{U} will give self-orthogonality after the following integer table is checked.

For the fixed family, the exact integer scalar products of incidence rows
are as follows. This table is checked for every pair of the $105$ flags by
the retained complete construction, not inferred from a sampled orbit.
\[
\begin{array}{l|r}
\text{relation between the two flags}&(RR^{\mathsf T})_{f,g}\\ \hline
f=g&24\\
\text{same point or same line, but distinct}&12\\
\text{distinct points and lines, with one cross incidence}&6\\
\text{distinct points and lines, with neither cross incidence}&3.
\end{array}
\]
Two distinct points and distinct lines cannot have both cross incidences.
Thus the table exhausts all cases and proves $RR^{\mathsf T}=0$ over
$\mathbb F_3$. Each row has sum $24$, so $Rj=0$, and the constant vector on
$120$ coordinates has norm zero. Set
\[
 K=\im(R^{\mathsf T}R),\qquad N=K+\mathbb F_3j.
\]
The exact row reductions give $\dim K=15$ and $\dim N=16$.
Consequently the sum defining $N$ is direct. Here the direct sum is an
internal decomposition of modules on the same $120$ coordinates, not a
concatenation of codes on disjoint coordinate sets. That common lemma proves
self-orthogonality of both codes. These ranks and the full table, including
its underlying block lists, are obligations of the complete finite certificate.

\subsection{Recognizing the socle}
The literal permutations in \texttt{module\_data.json} preserve the residual Steiner graph and generate a primitive
permutation group $G\cong L_3(4)$ on $\mathcal H$, with point stabilizer
$L_3(2)$. The certificate checks primitivity, the isomorphisms and stabilizer, so this is
the specified degree-$120$ coset action in the source. It also checks
invariance of $K$ and $N$. We apply the common full-matrix-algebra and socle criteria \cite[Theorems~6.2--6.3]{U}; the whole filtration and all intertwiner spaces must be checked.

Here $V=\mathbb F_3^{120}$. The supplied basis certifies invariant prefixes
of dimensions
\[
0,15,34,49,64,65,84,99,100,119,120.
\]
The ten factor dimensions and homomorphism-space dimensions are
\[
\begin{array}{c|rrrrrrrrrr}
i&1&2&3&4&5&6&7&8&9&10\\\hline
\dim D_i&15&19&15&15&1&19&15&1&19&1\\
\dim\operatorname{Hom}_G(D_i,V)&1&0&0&0&1&0&1&1&0&1.
\end{array}
\]
Triangularity of the generator matrices verifies invariance of all prefixes.
For each factor the matrix algebra has dimension $(\dim D_i)^2$, proving
simplicity by the common criterion. The first prefix is exactly $K$.
For each factor, the verifier forms \emph{all} equations
$XP_g=D_gX$ for the given generators and computes their full solution space.
Every row of every basis solution belongs to $N$. Linearity then establishes
the hypothesis for every homomorphism, not only a selected map.
The constant line is simple and $K$ is the first simple factor, so
the common criterion yields
\[
 \Soc(\mathbb F_3^{120})=N=K\oplus\mathbb F_3j.
\]
This calculation does not classify every irreducible module of $L_3(4)$.
The source's classification assertion concerning three nonisomorphic faithful
irreducible modules of dimension $15$ is retained from its reference [13], Jansen--Lux--Parker--Wilson, \emph{An Atlas of Brauer Characters} (1995).
Our certificate identifies a particular such simple module; its faithfulness
also follows because $G$ is simple and the representation is nontrivial.

\subsection{Weights and full permutation automorphisms}
Complete enumeration in the certified bases gives
\[
\begin{array}{c|r|r|r|r|r}
\text{code}&\text{number of words}&\text{dimension}&d&\text{words of weight }d&d^\perp\\\hline
K&14348907&15&48&1260&4\\
N&43046721&16&48&1260&4.
\end{array}
\]
The enumerator visits each of the $3^k$ coefficient vectors exactly once;
independence of the basis makes these distinct words. It records the full
weight distribution. The complete dual distributions follow from the common MacWilliams theorem \cite[Theorem~7.1]{U}. 
In both
cases the dual coefficients at weights $1,2,3$ vanish. At weight $4$ they are
$2100$ for $K^\perp$ and $1260$ for $N^\perp$. Thus the dual dimensions are $105$
and $104$, respectively, and the stated dual minimum distances follow.
These are complete finite computations, not distance bounds inferred from
self-orthogonality.

The orbit under $G$ of the supplied weight-$48$ word contains $1260$ distinct
words and $630$ distinct supports. Since $K\subset N$ and both codes have
exactly $1260$ minimum words, this orbit is the complete minimum-word set of
both codes. Every coordinate-permutation automorphism of either code
therefore acts on its support-incidence structure, with $120$ coordinate
vertices and $630$ distinct support vertices.

The common refinement theorem \cite[Theorem~6.4]{U} bounds the full automorphism group of this structure. 
Successively individualize the coordinate vertices
$1,6,48,10,2,3$ in the literal certificate ordering. Their successive
refined colour classes have sizes
\[
120,21,2,4,2,2.
\]
After these choices all coordinate classes are singletons. The theorem gives the upper bound $120\cdot21\cdot2\cdot4\cdot2\cdot2=80640$.

The explicit permutations supplied in the certificate generate a group
$A$ of order $80640$ and preserve each of $K,N$. Thus both code automorphism
groups attain the upper bound and equal $A$. The certificate checks
$A'=G$ and supplies two commuting involutions generating a complement of
order $4$, disjoint from $G$. Consequently
\[
 \Aut(K)=\Aut(N)=A\cong L_3(4):2^2.
\]
No graph-isomorphism program is needed for this upper bound. Such software
was useful in discovering the permutations, but the retained proof checks
literal permutations, ordinary group orders and the full refinement chain.

